\chapter{序論}
\label{chap:introduction}

\section{背景——出力の頑健性は，表現の頑健性を含意しない}
\label{sec:background}

\draftnote{context.md §1.1．VLM の頑健性は出力で測られてきたこと，Wani ら \cite{wani} が「答えが変わらないまま内部表現が大きく変化する」ことを示したこと（幾何的な加工と文字の重ね書きの違いまで），著者ら自身は対処法を示していないこと．}

\section{本研究の問い}
\label{sec:question}

\draftnote{context.md §1.2．Wani らは同じ入力の揺らぎを見ており，別々の質問どうしの近さは問われていない——そこから入る．中心の問いは一つ：タスク構造（署名）を内部表現に明示的に整合させたとき，内部表現と答えの生成に何が起きるか．正解率の向上は目的ではないことをここで明言する．}

\section{貢献}
\label{sec:contribution}

\draftnote{context.md §2．貢献は「読み取れること」と「使われること」を分けて示したこと．Elazar ら \cite{elazar} の「除く」方向に対し，本研究は損失で「足す」方向から構成的に示す．VLsI \cite{vlsi} は傍証にとどめ，中核の支柱には据えない（教師あり蒸留と自己整合の違い）．}

\section{本論文の構成}
\label{sec:organization}

\draftnote{各章の一文紹介．最後に書く．}
